% frontmatter/notation.tex
% Notation for the Parallel DFS chapter.

% Give every ppbox a blank line between its label and its body.
\tcbset{ppbox/.append style={after title={\par\smallskip}}}

\chapter*{Notation}
\addcontentsline{toc}{chapter}{Notation}

\begin{ppkey}
This page collects the symbols introduced in the Parallel Depth-First
Search chapter.  Entries are grouped by topic; within each group every
symbol is given its own definition box so the page can be skimmed by
eye when chasing down an unfamiliar piece of notation.
\end{ppkey}

%--------------------------------------------------------------
\section*{Asymptotic notation and basic mathematics}
%--------------------------------------------------------------

\begin{ppdefn}
\ppterm{Big-$\mathcal{O}$ notation} is an asymptotic upper bound:
$g(n) \in \mathcal{O}(f(n))$ if $|g(n)| \le c\,|f(n)|$ for some
constant $c$ and all sufficiently large $n$.
\end{ppdefn}

\begin{ppdefn}
\ppterm{Logarithm} $\log n$ denotes the base-2 logarithm by default.
\end{ppdefn}

\begin{ppdefn}
\ppterm{Set cardinality} $|S|$ is the number of elements in the finite
set $S$.  In particular $|C_i|$ is the number of vertices in
component $C_i$.
\end{ppdefn}

\begin{ppdefn}
\ppterm{Assignment} $\leftarrow$ means ``the variable on the left is
updated to the value on the right''.
\end{ppdefn}

%--------------------------------------------------------------
\section*{Graphs and trees}
%--------------------------------------------------------------

\begin{ppdefn}
\ppterm{Input graph} $G$ is a connected undirected graph; $n$ is the
number of vertices and $m$ the number of edges.
\end{ppdefn}

\begin{ppdefn}
\ppterm{Edge} $(u, v)$ is an undirected edge between vertices $u$ and
$v$.
\end{ppdefn}

\begin{ppdefn}
\ppterm{Directed arc} $u \to v$ is a directed arc from $u$ to $v$,
used for the arcs produced by an Euler tour of an undirected tree.
\end{ppdefn}

\begin{ppdefn}
\ppterm{Edge weight} $w$ is the weight of an edge, written inline as
$(u, v)\,w{=}k$ in worked examples.
\end{ppdefn}

\begin{ppdefn}
\ppterm{Root} $r$ is the root of a rooted spanning tree.
\end{ppdefn}

\begin{ppdefn}
\ppterm{Spanning tree} $S$ is a spanning tree of $G$, specifically the
(minimum) spanning tree built by the parallel Bor{\r u}vka rounds.
\end{ppdefn}

\begin{ppdefn}
\ppterm{DFS tree} $T$ is a DFS tree.  We write $T_i$ for the local
DFS tree of component $C_i$, and reserve ``the merged DFS tree'' for
the global tree assembled in Step~5.
\end{ppdefn}

\begin{ppdefn}
\ppterm{Parent pointer} $\texttt{parent[v]}$ is the parent of $v$ in
a rooted tree, treated as an array indexed by vertex id.
\end{ppdefn}

%--------------------------------------------------------------
\section*{Disjoint-set (Union-Find) operations}
%--------------------------------------------------------------

\begin{ppdefn}
\ppterm{Find} $\texttt{Find}(u)$ returns the representative (root) of
the set currently containing $u$.
\end{ppdefn}

\begin{ppdefn}
\ppterm{Union} $\texttt{Union}(u, v)$ merges the sets containing $u$
and $v$ into a single set.
\end{ppdefn}

\begin{ppdefn}
\ppterm{Pointer-jumping update} is the parallel assignment
$\texttt{parent[v]} \leftarrow \texttt{parent[parent[v]]}$.  Iterated
until a fixed point, it flattens any rooted forest in
$\mathcal{O}(\log n)$ rounds.
\end{ppdefn}

%--------------------------------------------------------------
\section*{Spine and off-spine components}
%--------------------------------------------------------------

\begin{ppdefn}
\ppterm{Spine} $P$ is the chosen root-to-leaf path in the spanning
tree, taken as the longest such path.  Written as a sequence
$s_0 \to s_1 \to \cdots \to s_d$.
\end{ppdefn}

\begin{ppdefn}
\ppterm{Spine vertex} $s_k$ is the spine vertex at position $k$ along
$P$.  $s_0$ is the root and $s_d$ the deepest leaf of the spanning
tree.
\end{ppdefn}

\begin{ppdefn}
\ppterm{Spine depth} $d$ is the depth of a spine vertex along $P$;
written next to a vertex as ``spine $d{=}k$''.
\end{ppdefn}

\begin{ppdefn}
\ppterm{Off-spine component} $C_i$ is the $i$-th connected component
of the graph induced when spine vertices are removed.
\end{ppdefn}

\begin{ppdefn}
\ppterm{Component root} $\mathrm{root}_i$ is the
deterministically-chosen entry vertex of $C_i$, used as the root of
the recursive DFS call into $C_i$.
\end{ppdefn}

\begin{ppdefn}
\ppterm{Off-spine parent} $\texttt{off\_parent[v]}$ is the auxiliary
parent pointer used to find off-spine components by parallel pointer
jumping.  At initialisation,
$\texttt{off\_parent[v]} = \texttt{parent[v]}$ if
$\texttt{parent[v]}$ is also off-spine, and
$\texttt{off\_parent[v]} = v$ otherwise.
\end{ppdefn}

%--------------------------------------------------------------
\section*{Euler tour and list ranking}
%--------------------------------------------------------------

\begin{ppdefn}
\ppterm{Next pointer} $\texttt{next}(u \to v)$ is the arc that follows
$u \to v$ in the Euler tour of the spanning tree, used as the
``next'' pointer for parallel list ranking.
\end{ppdefn}

\begin{ppdefn}
\ppterm{Down-arc} is an arc $u \to v$ where $v$ is a child of $u$ in
the rooted tree.  Down-arcs contribute $+1$ to the depth-counting
prefix sum.
\end{ppdefn}

\begin{ppdefn}
\ppterm{Up-arc} is an arc $v \to u$ where $u$ is the parent of $v$
(a return arc).  Up-arcs contribute $0$ to the depth-counting prefix
sum.
\end{ppdefn}

%--------------------------------------------------------------
\section*{Tree merging and edge classification}
%--------------------------------------------------------------

\begin{ppdefn}
\ppterm{Global DFS number} $\mathrm{dfs}[v]$ is the pre-order number
assigned to vertex $v$ in the merged DFS tree.  For a spine vertex
$s_k$:
$\mathrm{dfs}[s_k] = k + \sum_{j < k} \mathrm{block\_size}[s_j]$.
\end{ppdefn}

\begin{ppdefn}
\ppterm{Subtree size} $\mathrm{size}[v]$ is the number of vertices in
the subtree rooted at $v$ in the merged DFS tree.  Used together with
$\mathrm{dfs}[v]$ to test ancestor--descendant relations in
$\mathcal{O}(1)$.
\end{ppdefn}

\begin{ppdefn}
\ppterm{Block size} $\mathrm{block\_size}[s_j]$ is the size $|C_j|$
of the component hanging off spine vertex $s_j$ (zero if no component
hangs there).
\end{ppdefn}

\begin{ppdefn}
\ppterm{Tree edge} is an edge that belongs to the merged DFS tree.
\end{ppdefn}

\begin{ppdefn}
\ppterm{Back edge} is a non-tree edge $(u, v)$ such that one endpoint
is an ancestor of the other in the merged DFS tree.  In a valid DFS
tree of an undirected graph, every non-tree edge is a back edge.
\end{ppdefn}

\begin{ppdefn}
\ppterm{Forward edge} is, in a directed graph, a non-tree edge
$(u, v)$ where $v$ is a descendant of $u$ in the DFS tree.  Detected
by $\mathrm{dfs}[u] < \mathrm{dfs}[v]$ and
$\mathrm{dfs}[v] \le \mathrm{dfs}[u] + \mathrm{size}[u] - 1$.
\end{ppdefn}

\begin{ppdefn}
\ppterm{Cross edge} is a non-tree edge joining two vertices that are
neither ancestor nor descendant of each other.  In a valid undirected
DFS tree, no such non-tree edge exists.
\end{ppdefn}